\begin{definition}
Let (T) be an ordered tree-sequence, (v\in V(T)), (C\subseteq[k]),
and let (arphi,pi) be partial (k)-edge-colorings.  We say that
(pi) is ((T(v)-v,C,arphi))-stable if the following two conditions
hold on the finite vertex set (V(T(v)-v)) from
\noderef{TreeVerticesBefore}:
\begin{enumerate}
\item every edge identity incident with that set whose (arphi)-color
belongs to
\[
 \bigcup_{w\in V(T(v)-v)}\overline\varphi(w)\cup C
\]
has the same color under (pi) and (arphi);
\item (overline\pi(w)=\overline\varphi(w)) for every
(w\in V(T(v)-v)).
\end{enumerate}
This is the stability relation used by the paper for (T(v)-v), stated
directly on its vertices so that it remains meaningful when the preceding
display has no positive-length tree-sequence representation.
\end{definition}
